\sectioncentered*{Заключение}
\addcontentsline{toc}{section}{Заключение}

В рамках лабораторной работы №1 исследованы и программно реализованы механизмы избирательного управления доступом на двух уровнях: объектов файловой системы Linux и ролевой модели PostgreSQL. Экспериментальный прогон выполнен 10.09.2026.

\textbf{Семантическая составляющая.} Работа демонстрирует принципы ACL/DAC: в операционной системе~\mdash  через классы <<владелец / группа / остальные>> и биты \texttt{rwx}, в СУБД~\mdash  через роли и привилегии SQL. Показано, что эффективные права зависят не только от заданного набора битов или \texttt{GRANT}, но и от правил проверки доступа (в UNIX владелец не <<проваливается>> в групповые права) и от наследования ролей PostgreSQL.

\textbf{Прагматическая составляющая.} Подготовлен воспроизводимый стенд: автоматическое развёртывание, неразрушающая проверка политики и контролируемая очистка. Такой подход снижает риск ошибок ручной настройки и позволяет многократно демонстрировать и проверять ACL перед сдачей лабораторной.

Поставленные цели достигнуты: программы настройки и проверки доступа работоспособны; для PostgreSQL матрица операций дала \texttt{FAILS=0} (администратор~\mdash  полный доступ, пользователь~\mdash  CRUD без DDL, гость~\mdash  только SELECT).

\begin{thebibliography}{9}
	\addcontentsline{toc}{section}{Список использованных источников}

	\bibitem{olifer}
	Олифер В.\,Г., Олифер Н.\,А. Сетевые операционные системы: учебник для вузов. \mdash  2-е изд. \mdash  СПб. : Питер, 2009. \mdash  669~с.

	\bibitem{tanenbaum}
	Таненбаум Э.\,С. Современные операционные системы. \mdash  3-е изд. \mdash  СПб. : Питер, 2011. \mdash  1120~с.

	\bibitem{howard}
	Ховард М., Лебланк Д. Защищённый код. \mdash  М. : Русская редакция, 2004. \mdash  704~с.

	\bibitem{cwe}
	CWE: Common Weakness Enumeration [Электронный ресурс]. \mdash  Режим доступа: \url{http://cwe.mitre.org/}. \mdash  Дата доступа: 10.09.2026.

	\bibitem{postgres}
	PostgreSQL Documentation [Электронный ресурс]. \mdash  Режим доступа: \url{https://www.postgresql.org/docs/}. \mdash  Дата доступа: 10.09.2026.

	\bibitem{chmod}
	chmod(1) \mdash  Linux manual page [Электронный ресурс]. \mdash  Режим доступа: \url{https://man7.org/linux/man-pages/man1/chmod.1.html}. \mdash  Дата доступа: 10.09.2026.
\end{thebibliography}
